% ============ 2. Estado del arte (D2) ============
\section{Estado del arte y antecedentes}

Esta revisión sintetiza, \textbf{por enfoque} y no por autor, la literatura
relevante, distinguiendo lo hecho en el mundo de lo hecho en Colombia. Se apoya
en el estado del arte completo del equipo (20 págs., 39 referencias
verificadas).

\subsection{Enfoques 1--2: Ciencia de la Ciencia, evaluación responsable y sistemas nacionales}

La \emph{science of science} estudia el sistema científico con datos masivos y
métodos cuantitativos (Fortunato et al., 2018; Wang \& Barabási, 2021), y la
agenda de \textbf{evaluación responsable} —DORA (2012), Principios de Leiden
(Hicks et al., 2015), CoARA (2022)— fijó el estándar de uso de indicadores.
En América Latina, Brasil (Lattes) y Colombia (ScienTI-CVLAC/Publindex) son
referentes de registros curriculares (Revista CTS, 2010), con herramientas
como ScriptLattes (Mena-Chalco \& Cesar Júnior, 2009) y el método del
currículum vitae (Cañibano \& Bozeman, 2009). Para Colombia existen análisis
críticos del modelo de categorización ante el SNCTI (``Categorización de
grupos e investigadores'', 2023) y estudios de
métricas responsables y ciencia abierta (CONPES 4069, 2021). Nuestro proyecto
se inscribe en esta intersección: auditar un instrumento público de medición
bajo estándares de métricas responsables.

\subsection{Enfoques 3--4: calidad de datos y dinámica longitudinal}

La evaluación de registros administrativos usa las dimensiones de Wang \&
Strong (1996) y los principios FAIR (Wilkinson et al., 2016); la
reproducibilidad exige versionado de datos y observabilidad de pipelines
(Souza et al., 2023). El seguimiento de carreras se apoya en cadenas de Markov
(Kemeny \& Snell, 1976), modelos de supervivencia y análisis de cohortes; los
clásicos de física de carreras muestran que el impacto individual evoluciona
con reglas cuantificables (Sinatra et al., 2016; Petersen et al., 2012) y que
la movilidad geográfica estratifica la carrera (Deville et al., 2014; Sugimoto
et al., 2017).

\subsection{Enfoques 5--6: desigualdad, género, diversidad y redes}

La medición de la desigualdad usa HHI, Gini y Theil (Theil, 1967) y la
geografía de la ciencia (Frenken et al., 2009); la brecha de género global se
documentó en Larivière et al. (2013), con análisis interseccional (Crenshaw,
1989); las redes de colaboración (Newman, 2001) y la jerarquía institucional
extrema y persistente (Clauset et al., 2015) proveen los métodos de estructura
de poder del sistema.

\subsection{Contraste de aproximaciones y vacío}

Dos aproximaciones compiten para la pregunta longitudinal con supuestos
distintos: las \textbf{cadenas de Markov} suponen que la transición depende
solo del estado actual (homogeneidad), mientras los \textbf{modelos de
supervivencia} modelan el tiempo hasta el evento con covariables y censura.
Nuestro diseño \textbf{las contrasta}: primero se testea la propiedad
markoviana y luego se complementa con supervivencia. \textbf{Vacío que el
proyecto atiende:} no se ha publicado una auditoría estadística reproducible
del padrón de investigadores reconocidos de MinCiencias que combine, sobre los
datos abiertos oficiales, calidad del registro + trazabilidad longitudinal +
equidad con métodos de la ciencia de la ciencia; lo publicado en Colombia es
descriptivo o normativo, no evidencia cuantitativa auditable.
